\paragraph{Pattern checking.}
Checking a pattern against the remaining type (\ruleName{match-case}) rejects a tuple pattern at a type
with list values and a list pattern at a type with tuple values, and either at a type with range values, because matching either form of pattern
against the other kind of sequence or a range aborts with $\errExcluded$ (\secref{eval-patterns}). It also rejects a constructor
pattern at a type which does not determine the type arguments omitted by the pattern.

\paragraph{Shapes.}
\begin{figure}
\small
\added{\begin{tabular}[t]{rcll}
\grammarGroup{Heads} \\
$h$ & $::=$ & $\ell$ & (Literal) \\
    & $\mid$ & $C$ & (Class) \\
    & $\mid$ & $\headList{n}$ & (List of length $n$) \\
$H$ & $::=$ & $\set{h_1, \ldots, h_n}$ & (Excluded heads)
\end{tabular}
\hfill
\begin{tabular}[t]{rcll}
\grammarGroup{Shapes} \\
$k$ & $::=$ & $\shapeAny{\tau}{H}$ & (Rest) \\
    & $\mid$ & $\shapeAny{\patConstr{C[\seq{\tau}]}{k_1, \ldots, k_n}}{H}$ & (Constr) \\
    & $\mid$ & $\patTuple{k_1, \ldots, k_n}$ & (Tuple) \\
    & $\mid$ & $\shapeList{\tau}{k_1, \ldots, k_n}$ & (List) \\
    & $\mid$ & $\shapeDict{\tau}{\beta}{H}$ & (Dict) \\
$K$, $L$ & $::=$ & $\set{k_1, \ldots, k_n}$ & (Shape set)
\end{tabular}
}
\caption{\added{Heads and shapes}}
\label{fig:shapes}
\end{figure}
%
A shape (\figref{shapes}) denotes a set of values of a single type. A head is a literal, a class or a list
length, and denotes the values with that outermost form. The rest shape $\shapeAny{\tau}{H}$ denotes the
values of type $\tau$ whose head is not in the excluded set $H$. The other forms are constructor, tuple, list
and dictionary shapes, with shapes in place of sub-patterns; a constructor shape also has an excluded set,
containing the subclasses already matched. The auxiliary function $\shapes{\Sigma}(\tau, H)$ gives the
shapes of $\tau$ once the heads in $H$ are excluded.

\paragraph{Matching.}
The judgement $\Sigma; \Gamma \vdash L, p \matches K, L'$ matches a pattern $p$ against an ordered set of
shapes $L$, giving matched shapes $K$ and residual $L'$, and the judgement $\Sigma; \Gamma \vdash k, p
\matches K, L$ matches a pattern against a single shape $k$:
\begin{mathpar}
\small
\inferrule*[lab={\ruleName{pat-shapes}}]
{
  L' = \set{k \in L \mid \Sigma; \Gamma \vdash k, p \matches K_k, L_k} \\
  L' \neq \varnothing
}
{
  \Sigma; \Gamma \vdash L, p \matches \textstyle\bigdisjunion_{k \in L'} K_k, \; \textstyle\bigdisjunion_{k \in L'} L_k \disjunion (L \setminus L')
}

\and
\inferrule*[lab={\ruleName{pat-split}}]
{
  \Sigma; \Gamma \vdash k, p \splits K, L \\
  \Sigma; \Gamma \vdash K, p \matches K', L'
}
{
  \Sigma; \Gamma \vdash k, p \matches K', L \disjunion L'
}

\\
\inferrule*[lab={\ruleName{split-class}}]
{
  \Gamma \vdash q : \Type{C} \\
  \ancestors{\Sigma}(C) \cap H = \varnothing \\
  \Sigma \vdash C[\seq{\tau}'] \instanceAt \tau \\
  \meetTy{\Sigma}{\tau}{C[\seq{\tau}']} = D[\seq{\sigma}] \\
  \tyInst{\fields{\Sigma}(D)}{\seq{\sigma}} = \bind{\seq{x}}{\seq{\sigma}'} \\
  H' = \set{h \in H \mid \Sigma \vdash h : D[\seq{\sigma}]}
}
{
  \Sigma; \Gamma \vdash \shapeAny{\tau}{H}, \patConstr{q}{a} \splits \set{\shapeAny{\patConstr{D[\seq{\sigma}]}{\seq{k}}}{H'} \mid \seq{k} \in \shapes{\Sigma}(\seq{\sigma}')}, \shapes{\Sigma}(\tau, H \cup \set{C})
}

\\
\inferrule*[lab={\ruleName{pat-constr}}]
{
  \Gamma \vdash q : \Type{D} \\
  D \neq \clsRange \\
  D \in \ancestors{\Sigma}(C) \\
  \fields{\Sigma}(D) = \scheme{\seq{\alpha}}{\bind{\seq{x}}{\seq{\sigma}}} \\
  \fieldMap{\Sigma}(D, a) = \set{\bind{\seq{x}}{\seq{p}}} \\
  p'_i = \patWild \quad (\forall i) \\
  \Sigma; \Gamma \vdash \seq{k}, \seq{p} \mathbin{\concat} \seq{p}\,' \matches K, L
}
{
  \Sigma; \Gamma \vdash \shapeAny{\patConstr{C[\seq{\tau}]}{\seq{k}}}{H}, \patConstr{q}{a} \matches \set{\shapeAny{\patConstr{C[\seq{\tau}]}{\seq{k}\,'}}{H} \mid \seq{k}\,' \in K}, \set{\shapeAny{\patConstr{C[\seq{\tau}]}{\seq{l}}}{H} \mid \seq{l} \in L}
}

\end{mathpar}
The \ruleName{pat-shapes} rule matches the pattern against each shape of $L$. A case is reachable only if
its pattern matches at least one of the shapes, and the shapes which are not matched by the pattern pass into
the residual unchanged. A pattern matches a shape directly if the shape has the head tested by the pattern.
Otherwise \ruleName{pat-split} first splits the shape, by the judgement $\Sigma; \Gamma \vdash k, p \splits
K, L$, into shapes $K$ with that head and a residual $L$ with the head excluded. Splitting a rest shape at a
class $C$ (\ruleName{split-class}) gives constructor shapes for the instance of $C$ at the type of the shape,
with the shapes of the field types as components, and excludes $C$ from the residual; the instance
judgement, defined in \specfigref{instance}, recovers the type arguments omitted by the pattern. Shapes are split
only when a pattern requires it, because a class may refer to itself through its fields, and expanding every
shape in advance would not terminate. The \ruleName{pat-constr} rule matches a pattern for a class $D$
against a constructor shape for a subclass $C$ of $D$, matching the sub-patterns, with wildcards for omitted
fields, against the field shapes position by position.

The residual and the remaining type can differ. In the following program, the two cases between them match
every value of \pyinline{C}:

\begin{listingrow*}
\begin{listingcol}{0.45\textwidth}
\begin{python}[numbers=left, xleftmargin=1.5em]
@dataclass
class C:
    p: tuple[int, int] | None
\end{python}
\end{listingcol}
\begin{listingcol}{0.49\textwidth}
\begin{python}[numbers=left, firstnumber=4, xleftmargin=1.5em]
# error: f does not always return
def f(c: C) -> int:
    match c:
        case C(None):
            return 0
        case C((a, b)):
            return a + b
\end{python}
\end{listingcol}
\end{listingrow*}

\noindent The residual starts as the rest shape of \pyinline{C}. The first case splits it into a shape of
\pyinline{C} with a tuple field and a shape of \pyinline{C} with a \pyinline{None} field, and matches the
shape with the \pyinline{None} field. The second case matches the shape with the tuple field, leaving the
empty residual, so a third case would be unreachable.
Neither case matches every value of \pyinline{C}, however, so the remaining type is still \pyinline{C},
\ruleName{match-partial} applies, and the function is rejected as not always returning, following mypy. If
the second case were \pyinline{C(q)} instead, the remaining type after it would be $\tyNever$, but
\pyinline{q} would be bound at $\tyUnion{\tyTuple{\tyInt, \tyInt}}{\tyNone}$ rather than at $\tyTuple{\tyInt, \tyInt}$,
because pattern checking uses the remaining type.

